\documentclass{article}
\usepackage{amsmath, amssymb, amsthm}
\usepackage{hyperref}
\usepackage{geometry}
\geometry{margin=1in}

\title{Erd\H{o}s Problem \#383}
\date{Last updated: 2025-08-31}
\author{Source: erdosproblems.com}

\begin{document}
\maketitle

\noindent\textbf{Status:} OPEN \quad \textbf{No prize}

\noindent\textbf{Tags:} number theory

\medskip
\noindent\textbf{Problem Statement:}

Is it true that for every $k$ there are infinitely many primes $p$ such that the largest prime divisor of\[\prod_{0\leq i\leq k}(p^2+i)\]is $p$?

\bigskip
\noindent\textbf{Remarks:}

A positive answer to this would give an answer to the second part of [382]. Heuristically, the 'probability' that $n$ has no prime divisors $\geq n^{1/2}$ is $1-\log 2>0$, so standard heuristics predict the answer to this is yes.

\bigskip
\noindent\small{Source: \url{https://www.erdosproblems.com/383}}
\end{document}
